\documentclass[text06.tex]{subfiles}
\begin{document}
\newpage
\section{Open JTalkの声の\ruby{種類}{しゅ|るい}を変えてみる}
Open JTalkでは、声の種類や話し方を変えることができます。このような声の調子を\ruby{声音}{こわ|ね}といいます。普通の話し方、怒った話し方、うれしい話し方、悲しい話し方などから選ぶことができます。声の種類はhtsvoiceというファイルによって決められており、setvoice 命令で指定することができます。\\
\code{setvoice <声音ファイル>}

% TODO: relative or absolute path
<声音ファイル>には、 \ruby{拡張子}{かく|ちょう|し}が .htsvoice であるようなファイルへのパスを指定します。\\
みなさんの\mbox{Raspberry Pi}にはいくつかの声音ファイルがすでにダウンロードされています。  \textasciitilde /07/htsvoice-tohoku-f01を見てみてください。

\begin{itemize}
\item {\ttfamily htsvoice-tohoku-f01/\allowbreak tohoku-f01-neutral.htsvoice}（普通の話し方）
\item {\ttfamily htsvoice-tohoku-f01/\allowbreak tohoku-f01-angry.htsvoice}（怒った話し方）
\item {\ttfamily htsvoice-tohoku-f01/\allowbreak tohoku-f01-happy.htsvoice}（うれしい話し方）
\item {\ttfamily htsvoice-tohoku-f01/\allowbreak tohoku-f01-sad.htsvoice}（悲しい話し方）
\end{itemize}

が利用できます。例を示します。 \textasciitilde /07/voice.hsp をHSPスクリプトエディタで開いてみましょう。\\

\begin{lstlisting}[caption=voice.hsp,label=voice.hsp]
#include "hsp3dish.as"
#include "jtalk.as"

redraw 0
mes "発話開始"
mes "「子どもIT未来塾」"
redraw 1

setvoice "htsvoice-tohoku-f01/tohoku-f01-sad.htsvoice"	<#blue#;音声ファイル指定#>
jtload "子どもIT未来塾", 0				<#blue#;音声ファイル生成と読み込み#>
mmplay 0						<#blue#;音声ファイルの再生#>
stop
\end{lstlisting}

8行目までは前と同じです。9行目でsetvoiceが使用されています\\
\code{setvoice "htsvoice-tohoku-f01/\allowbreak tohoku-f01-sad.htsvoice"}\\
\code{voice.hsp} は \code{\textasciitilde /07} にあるので、
\code{voice.hsp} 内に書かれた {\ttfamily htsvoice-tohoku-f01/\allowbreak tohoku-f01-sad.htsvoice}
は \code{\textasciitilde /07} から\ruby{探}{さが}されます。
この行\ruby{以降}{い|こう}、jtloadが\ruby{呼}{よ}び出されると {\ttfamily \textasciitilde /07/htsvoice-tohoku-f01/\allowbreak tohoku-f01-sad.htsvoice}
という声音が使用されます。ダブルクォーテーションで囲まれている部分を\ruby{変更}{へん|こう}することで、別の声の種類にすることができます。\\

\begin{tcolorbox}[title=\useOmetoi]
\begin{enumerate}
\addex{voice.hsp をHSPスクリプトエディタで実行してみましょう。}
\addex{時間があったらやりましょう。\\好きな声音を使って読み上げさせましょう。}
\end{enumerate}
\end{tcolorbox}
\end{document}
